% ======================================================================
% Chapter 7: Radiative Corrections: Some Formal Developments
% (book pages 211-257; PDF pages 234-287)
% ======================================================================
\chapter{辐射修正：若干形式发展}

在过去的三个章节中，我们曾经四次“作弊”，%
\footnote{第五次“作弊”是直接假定而不是推导光子传播子，这将在第9章中得到补救。}即陈述（有时是给出动机）一个结果，但推迟了它的证明。这些结果分别是：
\begin{enumerate}
  \item 用 $S$ 矩阵元表示衰变率的公式，式~(4.86)。
  \item 用费曼图表示 $S$ 矩阵元的主公式，式~\eqref{eq:4.103}。
  \item 沃德恒等式，式~\eqref{eq:5.79}。
  \item 为消除顶点修正图中的紫外发散而进行的特别减除，式~(6.55)。
\end{enumerate}
现在是时候回到这些问题，对它们进行恰当的处理了。在第7.2节至第7.4节中，我们将推导这全部四个结果。在此过程中获得的知识将帮助我们解读 \eqref{eq:6.1} 所示的电子在重靶上散射的其余三个圈修正：外腿修正和真空极化。我们将在第7.1节计算前者，在第7.5节计算后者。

本章将比前两章更加抽象。它的主要主题是作为外动量解析函数的费曼图的奇点。然而，我们将发现，这个表面上深奥的主题蕴含着丰富的物理内涵，并阐明了费曼图与量子理论一般原理之间的关系。

\section{场强重整化}

在本节中，我们将研究两点关联函数
\begin{equation}
  \langle\Omega|T\phi(x)\phi(y)|\Omega\rangle
  \qquad\text{或}\qquad
  \langle\Omega|T\psi_\alpha(x)\bar{\psi}_\beta(y)|\Omega\rangle .
\end{equation}
的解析结构。在自由场论中，两点函数
$\langle 0|T\phi(x)\phi(y)|0\rangle$ 有一个简单的解释：它是粒子从 $y$ 传播到 $x$ 的振幅。这种解释在多大程度上可以推广到相互作用理论呢？

我们对两点函数的分析将只依赖于相对论和量子力学的一般原理；它不依赖于相互作用的性质，也不依赖于微扰论的展开。不过，我们将把讨论限制在标量场。对于带自旋场的关联函数，可以得到类似的结果；我们将在分析结束时展示 Dirac 场的类似结果。

为了剖析两点函数 $\langle\Omega|T\phi(x)\phi(y)|\Omega\rangle$，我们将在 $\phi(x)$ 与 $\phi(y)$ 之间插入单位算符，其形式为对一组完备态的求和。我们选取这些态为完整的相互作用哈密顿量 $H$ 的本征态。由于动量算符 $\mathbf{P}$ 与 $H$ 对易，我们也可以选取这些态为 $\mathbf{P}$ 的本征态。但我们还可以更充分地利用洛伦兹不变性。设 $|\lambda_0\rangle$ 是 $H$ 的动量为零的本征态：$\mathbf{P}|\lambda_0\rangle = 0$。则 $|\lambda_0\rangle$ 的所有加速（boost）也都是 $H$ 的本征态，并且它们具有一切可能的三动量。反之，$H$ 的任何具有确定动量的本征态都可以写成某个零动量本征态 $|\lambda_0\rangle$ 的加速。四动量算符 $P^\mu = (H, \mathbf{P})$ 的本征值排列成一系列双曲面，如图~\ref{fig:7.1} 所示。

\begin{figure}[t]
  \centering
  \includegraphics[width=0.55\textwidth]{fig7.1.png}
  \caption{四动量算符 $P^\mu = (H, \mathbf{P})$ 的本征值在能量-动量空间中占据一系列双曲面。对于典型理论，态由一个或多个质量为 $m$ 的粒子组成。因此存在一个单粒子态的双曲面，以及两粒子态、三粒子态等的双曲面连续统。在产生两个自由粒子的阈值之下，还可能存在一个或多个束缚态双曲面。}
  \label{fig:7.1}
\end{figure}

回忆第2章，单粒子态的完备性关系为
\begin{equation}
  (\text{恒等算符})_{\text{单粒子}} = \int
    \frac{d^3p}{(2\pi)^3}\ \frac{1}{2E_{\mathbf{p}}}\ |\mathbf{p}\rangle
    \langle\mathbf{p}| .
  \label{eq:7.1}
\end{equation}
借助一些记号，我们可以为整个 Hilbert 空间写出类似的完备性关系。设 $|\lambda_{\mathbf{p}}\rangle$ 是 $|\lambda_0\rangle$ 的具有动量 $\mathbf{p}$ 的加速，并假设态 $|\lambda_{\mathbf{p}}\rangle$ 与单粒子态 $|\mathbf{p}\rangle$ 一样按相对论方式归一化。设
$E_{\mathbf{p}}(\lambda) = \sqrt{|\mathbf{p}|^2 + m_\lambda^2}$，其中
$m_\lambda$ 是态 $|\lambda_{\mathbf{p}}\rangle$ 的“质量”，即态 $|\lambda_0\rangle$ 的能量。于是所需的完备性关系为
\begin{equation}
  1 = |\Omega\rangle\langle\Omega| + \sum_\lambda \int
    \frac{d^3p}{(2\pi)^3}\ \frac{1}{2E_{\mathbf{p}}(\lambda)}\
    |\lambda_{\mathbf{p}}\rangle\langle\lambda_{\mathbf{p}}|,
  \label{eq:7.2}
\end{equation}
其中求和遍历所有零动量态 $|\lambda_0\rangle$。

现在我们在两点函数的算符之间插入这一展开。暂时假设 $x^0 > y^0$。让我们略去无趣的常数项
$\langle\Omega|\phi(x)|\Omega\rangle\langle\Omega|\phi(y)|\Omega\rangle$。
（这一项通常由对称性而为零；对于高自旋场，它由洛伦兹不变性而为零。）于是两点函数为
\begin{equation}
  \langle\Omega|\phi(x)\phi(y)|\Omega\rangle
  = \sum_\lambda \int \frac{d^3p}{(2\pi)^3}\ \frac{1}{2E_{\mathbf{p}}(\lambda)}\
    \langle\Omega|\phi(x)|\lambda_{\mathbf{p}}\rangle\
    \langle\lambda_{\mathbf{p}}|\phi(y)|\Omega\rangle .
  \label{eq:7.3}
\end{equation}
我们可以如下处理这些矩阵元：
\begin{equation}
  \begin{aligned}
    \langle\Omega|\phi(x)|\lambda_{\mathbf{p}}\rangle
    &= \langle\Omega|e^{iP\cdot x}\phi(0)e^{-iP\cdot x}|\lambda_{\mathbf{p}}\rangle\\
    &= \langle\Omega|\phi(0)|\lambda_{\mathbf{p}}\rangle\ e^{-ip\cdot x}
      \Big|_{p^0=E_{\mathbf{p}}}\\
    &= \langle\Omega|\phi(0)|\lambda_0\rangle\ e^{-ip\cdot x}
      \Big|_{p^0=E_{\mathbf{p}}} .
  \end{aligned}
  \label{eq:7.4}
\end{equation}
最后一个等式是 $|\Omega\rangle$ 和 $\phi(0)$ 的洛伦兹不变性的结果：插入 $U$ 因子，其中 $U$ 是实现从 $\mathbf{p}$ 到 0 的洛伦兹加速的幺正算符，并利用
$U\phi(0)U^{-1} = \phi(0)$。（对于带自旋的场，我们需要跟踪其非平凡的洛伦兹变换。）引入对 $p^0$ 的积分，我们关于两点函数的表达式（仍假设 $x^0 > y^0$）变为
\begin{equation}
  \langle\Omega|\phi(x)\phi(y)|\Omega\rangle
  = \sum_\lambda \int \frac{d^4p}{(2\pi)^4}\ \frac{i}{p^2-m_\lambda^2+i\epsilon}\
    e^{-ip\cdot(x-y)}\ \bigl| \langle\Omega|\phi(0)|\lambda_0\rangle
      \bigr|^2 .
  \label{eq:7.5}
\end{equation}
注意这里出现了费曼传播子 $D_F(x-y)$，只是把 $m$ 替换成了 $m_\lambda$。

\begin{figure}[t]
  \centering
  \includegraphics[width=0.5\textwidth]{fig7.2.png}
  \caption{典型相互作用场论的谱函数 $\rho(M^2)$。单粒子态在 $m^2$（粒子质量的平方）处贡献一个 delta 函数。多粒子态具有从 $(2m)^2$ 开始的连续谱。也可能存在束缚态。}
  \label{fig:7.2}
\end{figure}

对于 $y^0 > x^0$ 的情形，有类似的表达式。这两种情形都可以概括为两点函数的如下一般表示（K\"all\'en--Lehmann 谱表示）：
\begin{equation}
  \langle\Omega|T\phi(x)\phi(y)|\Omega\rangle
  = \int_0^{\infty} \frac{dM^2}{2\pi}\ \rho(M^2)\ D_F(x-y; M^2),
  \label{eq:7.6}
\end{equation}
其中 $\rho(M^2)$ 是正的谱密度函数，
\begin{equation}
  \rho(M^2) = \sum_\lambda (2\pi)\,\delta(M^2 - m_\lambda^2)\,
    \bigl| \langle\Omega|\phi(0)|\lambda_0\rangle \bigr|^2 .
  \label{eq:7.7}
\end{equation}
典型理论的谱密度 $\rho(M^2)$ 画在图~\ref{fig:7.2} 中。注意单粒子态对谱密度贡献一个孤立的 delta 函数：
\begin{equation}
  \rho(M^2) = 2\pi\,\delta(M^2-m^2)\cdot Z
    + (\text{在 } M^2 \ge (2m)^2 \text{ 之前没有其他贡献}),
  \label{eq:7.8}
\end{equation}
其中 $Z$ 是由 \eqref{eq:7.7} 中的矩阵元平方给出的某个数。我们称 $Z$ 为\emph{场强重整化}。量 $m$ 是单个粒子的精确质量——静止时能量的精确本征值。一般而言，这个量与出现在拉格朗日量中的质量参数值不同。我们将把拉格朗日量中的参数记作 $m_0$，即裸质量，并把 $m$ 称为 $\phi$ 玻色子的物理质量。只有物理质量 $m$ 是直接可观测的。

谱分解 \eqref{eq:7.6} 给出两点函数 Fourier 变换的如下形式：
\begin{equation}
  \begin{aligned}
    \int d^4x\ \langle\Omega|T\phi(x)\phi(0)|\Omega\rangle\ e^{ip\cdot x}
    &= \frac{iZ}{p^2-m^2+i\epsilon}
      + \int_{4m^2}^{\infty} \frac{dM^2}{2\pi}\ \rho(M^2)\
        \frac{i}{p^2-M^2+i\epsilon} .
  \end{aligned}
  \label{eq:7.9}
\end{equation}
这个函数在复 $p^2$ 平面上的解析结构如图~\ref{fig:7.3} 所示。第一项在 $p^2 = m^2$ 处给出一个孤立的单极点，而第二项给出从两粒子阈值 $p^2 = 4m^2$ 开始的支割线。

\begin{figure}[t]
  \centering
  \includegraphics[width=0.5\textwidth]{fig7.3.png}
  \caption{典型理论的两点函数 Fourier 变换在复 $p^2$ 平面上的解析结构。单粒子态在粒子质量的平方处贡献一个孤立极点。两个或更多自由粒子的态给出支割线，而束缚态给出额外的极点。}
  \label{fig:7.3}
\end{figure}

对于自旋 $1/2$ 的场，\eqref{eq:7.9} 的类似式为
\begin{equation}
  \int d^4x\ \langle\Omega|T\psi_\alpha(x)\bar{\psi}_\beta(0)|\Omega\rangle\
    e^{ip\cdot x}
  = \frac{i(\slashed{p}+m)\,Z_2}{p^2-m^2+i\epsilon}
    + \int_{4m^2}^{\infty} \frac{dM^2}{2\pi}\ \rho(M^2)\
      \frac{i(\slashed{p}+M)}{p^2-M^2+i\epsilon},
  \label{eq:7.12}
\end{equation}
其中 $Z_2$ 是 Dirac 场的场强重整化，$\rho(M^2)$ 是相应的谱函数。对于更高自旋的场，\eqref{eq:7.9} 的类似式通过乘以适当的极化张量得到。

\subsection{电子自能}

现在让我们在微扰论中计算电子传播子的领头修正，并确定场强重整化常数 $Z_2$ 和质量移动 $\delta m = m - m_0$。我们在 QED 中工作。对电子自能贡献的最低阶图为
\begin{equation}
  \text{（单圈电子自能图）} .
  \label{eq:7.15}
\end{equation}
利用费曼规则，其值为
\begin{equation}
  \begin{aligned}
    \Sigma_2(\slashed{p}) &= (-ie)^2\int \frac{d^4k}{(2\pi)^4}\
      \gamma^\mu\ \frac{i(\slashed{k}+m_0)}{k^2-m_0^2+i\epsilon}\
      \gamma^\nu\ \frac{-ig_{\mu\nu}}{(p-k)^2-\mu^2+i\epsilon}\\
    &= -e^2\int \frac{d^4k}{(2\pi)^4}\
      \frac{\gamma^\mu(\slashed{k}+m_0)\gamma_\mu}
        {[k^2-m_0^2+i\epsilon][(p-k)^2-\mu^2+i\epsilon]} .
  \end{aligned}
  \label{eq:7.16}
\end{equation}
利用收缩恒等式
\begin{equation}
  \gamma^\mu\gamma_\mu = 4,\qquad
  \gamma^\mu\slashed{k}\gamma_\mu = -2\slashed{k},
\end{equation}
分子简化为
\begin{equation}
  \gamma^\mu(\slashed{k}+m_0)\gamma_\mu = 2m_0 - 2\slashed{k} .
\end{equation}
引入费曼参数来合并分母，
\begin{equation}
  \frac{1}{[k^2-m_0^2][(p-k)^2-\mu^2]}
  = \int_0^1 dx\ \frac{1}{[(k-xp)^2 - xp^2 + xp^2 + (1-x)m_0^2 + x\mu^2]^2}
  \cdots ,
\end{equation}
并作平移 $\ell = k - xp$，我们得到
\begin{equation}
  \Sigma_2(\slashed{p}) = 2e^2\int_0^1 dx\ \int
    \frac{d^4\ell}{(2\pi)^4}\ \frac{-2x\slashed{p} + 4m_0}
      {[\ell^2 - \Delta + i\epsilon]^2},
  \label{eq:7.17}
\end{equation}
其中 $\Delta = -x(1-x)p^2 + x\mu^2 + (1-x)m_0^2$。对 $\ell$ 的积分是发散的。为了计算它，我们先用 Pauli--Villars 程序对之正则化：
\begin{equation}
  \frac{1}{(p-k)^2-\mu^2+i\epsilon}
    \to \frac{1}{(p-k)^2-\mu^2+i\epsilon}
      - \frac{1}{(p-k)^2-\Lambda^2+i\epsilon} .
\end{equation}
第二项将与 \eqref{eq:7.17} 具有相同的形式，只是把 $\mu$ 替换为 $\Lambda$。与第6.3节一样，我们现在作 Wick 转动并代换欧几里得变量 $\ell_E^0 = -i\ell^0$。这给出
\begin{equation}
  \int \frac{d^4\ell}{(2\pi)^4}\ \frac{1}{[\ell^2-\Delta]^2}
  = \frac{i}{(4\pi)^2}\int_0^{\infty} d\ell_E\
    \frac{\ell_E^3}{[\ell_E^2+\Delta]^2}
  = \frac{i}{(4\pi)^2}\ \log\Bigl( \frac{\Lambda^2}{\Delta} \Bigr),
  \label{eq:7.18}
\end{equation}
其中 $\Delta_\Lambda = -x(1-x)p^2 + x\Lambda^2 + (1-x)m_0^2 \to x\Lambda^2$
当 $\Lambda\to\infty$ 时。因此最终结果为
\begin{equation}
  \Sigma_2(\slashed{p}) = \frac{\alpha}{2\pi}\int_0^1 dx\ (2m_0 - x\slashed{p})\
    \log\Bigl( \frac{x\Lambda^2}{(1-x)m_0^2 + x\mu^2 - x(1-x)p^2} \Bigr) .
  \label{eq:7.19}
\end{equation}
在讨论这个表达式中的发散之前，让我们先研究它作为 $p^2$ 函数的解析行为。\eqref{eq:7.19} 中的对数当其宗量为负时具有支割线，并且对于任何固定的 $x$，这将在 $p^2$ 足够大时发生。更确切地说，割线起始于
\begin{equation}
  (1-x)m_0^2 + x\mu^2 - x(1-x)p^2 = 0 .
\end{equation}
解出这个方程中的 $x$，我们发现 $\Sigma_2(p^2)$ 的支割线起始于 $p^2 = (m_0+\mu)^2$，即在两粒子（电子加光子）态产生的阈值处。%
\footnote{在真实的 QED 中，$\mu = 0$，两粒子支割线与单粒子极点合并。这个微妙之处在红外发散的抵消的完整处理中起作用，但超出了我们当前分析的范畴。}支割线的位置恰好是我们在 K\"all\'en--Lehmann 谱表示中所预期的位置。

我们现在已经找到了 K\"all\'en--Lehmann 表示所预言的两粒子支割线，但还没有找到在 $p^2 = m^2$ 处所预期的简单极点。为了找到它，我们必须真正计入无穷级数的费曼图。幸运的是，这个级数很容易求和。

让我们定义一个\emph{单粒子不可约}（1PI）图为任何不能通过移去一条线而分成两部分的图：
\begin{equation}
  \text{（具有自能插入的图是 1PI）} .
\end{equation}
设 $-i\Sigma(\slashed{p})$ 表示所有具有两条外费米子线的 1PI 图之和：
\begin{equation}
  -i\Sigma(\slashed{p}) = \text{（所有 1PI 两点图之和）} .
  \label{eq:7.21}
\end{equation}
（尽管每个图有两条外线，但这些线的费曼传播子不包括在 $\Sigma(p)$ 的表达式中。）在 $\alpha$ 的最低阶，我们看到 $\Sigma = \Sigma_2$。

两点函数的 Fourier 变换现在可以写成
\begin{equation}
  \int d^4x\ \langle\Omega|T\psi(x)\bar{\psi}(0)|\Omega\rangle\ e^{ip\cdot x}
  = \frac{i}{\slashed{p}-m_0}
    + \frac{i}{\slashed{p}-m_0}\ (-i\Sigma)\ \frac{i}{\slashed{p}-m_0}
    + \cdots .
  \label{eq:7.22}
\end{equation}
第一个图在 $p^2 = m_0^2$ 处有一个简单极点。第二类中的每个图在 $p^2 = m_0^2$ 处有一个二重极点。第三类中的每个图有一个三重极点。当我们计入越来越多的图时，$p^2 = m_0^2$ 附近的行为越来越糟。但幸运的是，所有图之和构成一个几何级数。注意 $\Sigma(\slashed{p})$ 与 $\slashed{p}$ 对易，因为 $\Sigma(\slashed{p})$ 只是纯数和 $\slashed{p}$ 的函数。事实上，我们可以把 $\Sigma(\slashed{p})$ 看成 $\slashed{p}$ 的函数，写出 $p^2 = (\slashed{p})^2$。于是我们可以把每个电子传播子改写为 $i/(\slashed{p}-m_0)$，并把上述级数表示为
\begin{equation}
  \begin{aligned}
    \int d^4x\ \langle\Omega|T\psi(x)\bar{\psi}(0)|\Omega\rangle\ e^{ip\cdot x}
    &= \frac{i}{\slashed{p}-m_0}\ \Bigl[ 1 +
      \frac{\Sigma(\slashed{p})}{\slashed{p}-m_0}
      + \Bigl( \frac{\Sigma(\slashed{p})}{\slashed{p}-m_0} \Bigr)^2
      + \cdots \Bigr]\\
    &= \frac{i}{\slashed{p}-m_0-\Sigma(\slashed{p})} .
  \end{aligned}
  \label{eq:7.23}
\end{equation}
完整传播子有一个简单极点，它被 $\Sigma(\slashed{p})$ 从 $m_0$ 移动开了。这个极点的位置，即物理质量 $m$，是方程
\begin{equation}
  \bigl[ \slashed{p} - m_0 - \Sigma(\slashed{p}) \bigr]
    \Big|_{\slashed{p}=m} = 0 .
  \label{eq:7.24}
\end{equation}
的解。注意，如果 $\Sigma(\slashed{p})$ 由约定 \eqref{eq:7.21} 定义，那么对 $\Sigma$ 的正贡献给出电子质量的正移动。在极点附近，\eqref{eq:7.23} 的分母具有形式
\begin{equation}
  (\slashed{p}-m)\cdot \Bigl( 1 - \frac{d\Sigma}{d\slashed{p}}
    \Big|_{\slashed{p}=m} \Bigr) + \mathcal{O}\bigl( (\slashed{p}-m)^2 \bigr)
  = (\slashed{p}-m)\cdot Z_2^{-1} + \cdots .
  \label{eq:7.25}
\end{equation}
因此完整的电子传播子具有恰好 \eqref{eq:7.12} 形式的单粒子极点，其中 $m$ 由 \eqref{eq:7.24} 给出，且
\begin{equation}
  Z_2^{-1} = 1 - \frac{d\Sigma}{d\slashed{p}}\Big|_{\slashed{p}=m} .
  \label{eq:7.26}
\end{equation}
我们对 $\Sigma_2$ 的显式计算使我们能够计算 $m$ 和 $Z_2$ 的第一修正。让我们从 $m$ 开始。在 $\alpha$ 阶，质量移动为
\begin{equation}
  \delta m = m - m_0 = \Sigma_2(\slashed{p}=m) \approx \Sigma_2(\slashed{p}=m_0) .
  \label{eq:7.27}
\end{equation}
完成积分，人们发现 $\delta m$ 是对数发散（紫外）的，并且还包含一个红外对数。紫外发散通过裸质量的重整化被吸收，而红外发散与轫致辐射相抵消，正如我们在第6章所讨论的那样。类似地，$\Sigma_2$ 在 $\slashed{p} = m$ 处的导数确定 $Z_2$；在 $\alpha$ 阶，
\begin{equation}
  Z_2 = 1 - \frac{d\Sigma_2}{d\slashed{p}}\Big|_{\slashed{p}=m}
    + \mathcal{O}(\alpha^2) .
  \label{eq:7.29}
\end{equation}

\section{LSZ约化公式}

我们现在想把理论中的 $S$ 矩阵元——例如 $2\to2$ 过程的散射振幅——表示为场的时间编序乘积的真空期望值的 Fourier 变换，即关联函数。这个关系由 LSZ
(Lehmann--Symanzik--Zimmermann) 约化公式提供，我们在本节中推导它。%

\subsection*{LSZ 公式的推导}

考虑一个具有确定动量 $\mathbf{p}$ 的单粒子态在较晚时刻演化到同一态的振幅。在相互作用理论中，在遥远过去自由演化的态 $|\mathbf{p}\rangle$ 与完整哈密顿量的具有确定动量的本征态并不相同。相反，出现在谱分解中的态是物理单粒子态，其归一化为
\begin{equation}
  \langle\mathbf{p}|\mathbf{q}\rangle = 2E_{\mathbf{p}}\,(2\pi)^3\,
    \delta^{(3)}(\mathbf{p}-\mathbf{q}) .
\end{equation}
场 $\phi(x)$ 在真空与这个物理单粒子态之间具有非零矩阵元，
\begin{equation}
  \langle\Omega|\phi(x)|\mathbf{p}\rangle = \sqrt{Z}\,
    e^{-ip\cdot x},
  \label{eq:7.35}
\end{equation}
其中 $Z$ 是式~\eqref{eq:7.8} 的场强重整化常数。

为了推导约化公式，考虑具有两个外粒子的关联函数：
\begin{equation}
  \langle\Omega|T\{\phi(x_1)\phi(x_2)\cdots\}|\Omega\rangle .
  \label{eq:7.41}
\end{equation}
在诸场之间插入一组完备态，并通过把坐标分离成粒子充分分离的区域（遥远的过去和遥远的未来）以及它们相互作用的有界区域来分析时间积分。这一分析的结果是，当波包取为确定动量时，能量变量中的领头奇点给出 $S$ 矩阵元。通过交换极限的顺序，得到 LSZ 约化公式：
\begin{equation}
  \begin{aligned}
    &\Bigl[ \prod_{i=1}^{n} \int d^4x_i\ e^{ip_i\cdot x_i}\
      \frac{i}{p_i^2-m^2+i\epsilon}\ \sqrt{Z}\ \Bigr]\
    \Bigl[ \prod_{j=1}^{m} \int d^4y_j\ e^{-ik_j\cdot y_j}\
      \frac{i}{k_j^2-m^2+i\epsilon}\ \sqrt{Z}\ \Bigr]\\
    &\quad \times \langle\Omega|T\{\phi(x_1)\cdots\phi(x_n)\phi(y_1)\cdots
      \phi(y_m)\}|\Omega\rangle\\
    &\xrightarrow{\text{每个 } p_i^0\to +E_{\mathbf{p}_i};\ \text{每个 }
      k_j^0\to +E_{\mathbf{k}_j}}
    \langle\mathbf{p}_1\cdots\mathbf{p}_n|S|\mathbf{k}_1\cdots
      \mathbf{k}_m\rangle .
  \end{aligned}
  \label{eq:7.42}
\end{equation}
出现在这个方程中的量 $Z$ 恰好是场强重整化常数，在第7.1节中被定义为场的两点函数中单粒子极点的留数。每个不同的粒子都将有各自的重整化因子 $Z$，从它自己的两点函数得到。对于高自旋场，每个 $\sqrt{Z}$ 因子都带有一个极化因子，如 $u^s(p)$，正如式~\eqref{eq:7.12} 中那样。在 \eqref{eq:7.42} 的第二行中，必须对极化 $s$ 求和。

用文字来说，LSZ 公式说一个 $S$ 矩阵元可以如下计算。计算适当的 Fourier 变换后的关联函数，看外粒子靠近质壳的动量空间区域，并识别多粒子极点的系数。对于带自旋的场，还必须乘以一个极化旋量（如 $u^s(p)$）或矢量（如 $\epsilon^r(k)$）以投影出所需的自旋态。

我们的下一个目标是用费曼图的语言来表达这一程序。让我们分析标量场四点函数的图解展开与 2 粒子 $\to$ 2 粒子散射的 $S$ 矩阵元之间的关系。我们将显式地考虑对关联函数有贡献的全连通费曼图。通过类似的分析，很容易确认非连通图应当被舍弃，因为它们不具有 \eqref{eq:7.42} 右边所示的四个极点乘积的奇点结构。

精确的四点函数
\begin{equation}
  \Bigl( \prod_{i=1}^{2}\int d^4x_i\ e^{ip_i\cdot x_i}\ \Bigr)\
    \Bigl( \prod_{j=1}^{2}\int d^4y_j\ e^{-ik_j\cdot y_j}\ \Bigr)\
    \langle\Omega|T\{\phi(x_1)\phi(x_2)\phi(y_1)\phi(y_2)\}|\Omega\rangle
\end{equation}
具有图~\ref{fig:7.4} 所示的一般形式。在这个图中，我们显式地标出了每条腿上的图解修正；中心阴影圆圈代表所有切掉外腿（amputated）的四点图之和。

\begin{figure}[t]
  \centering
  \includegraphics[width=0.55\textwidth]{fig7.4.png}
  \caption{标量场论中精确四点函数的结构。外腿携带场强重整化因子 $\sqrt{Z}$，阴影圆圈是所有切掉外腿的 1PI 四点图之和。}
  \label{fig:7.4}
\end{figure}

我们可以像上一节中对电子传播子所做的那样，对每条外腿的修正求和。设 $-iM^2(p^2)$ 表示对标量传播子的所有单粒子不可约（1PI）插入之和：
\begin{equation}
  -iM^2(p^2) = \text{（1PI 自能插入之和）} .
\end{equation}
于是精确传播子可以写成几何级数并按式~\eqref{eq:7.23} 求和：
\begin{equation}
  \frac{i}{p^2-m_0^2-M^2(p^2)} .
  \label{eq:7.43}
\end{equation}
注意，与电子传播子的情形一样，我们对 1PI 自能 $M^2(p^2)$ 的符号约定意味着 $M^2(p^2)$ 的正贡献对应于标量粒子质量的正移动。如果我们把每个重求和后的传播子围绕物理粒子极点展开，我们看到四点振幅的每条外腿贡献
\begin{equation}
  \frac{iZ}{p^2-m^2-M^2} \xrightarrow{p^0\to E_{\mathbf{p}}}
    \frac{iZ}{p^2-m^2} + (\text{正则部分}) .
  \label{eq:7.44}
\end{equation}
因此，图之和包含四个极点的乘积：
\begin{equation}
  \frac{iZ}{p_1^2-m^2}\ \frac{iZ}{p_2^2-m^2}\ \frac{iZ}{k_1^2-m^2}\
    \frac{iZ}{k_2^2-m^2}\ \times\ \text{（切掉外腿的连通振幅）} .
\end{equation}
LSZ 公式 \eqref{eq:7.42} 指示我们把关联函数乘以
$\prod (i\sqrt{Z}/(p^2-m^2+i\epsilon))$，然后在极点处取留数。因子 $i/(p^2-m^2)$ 与外腿上的传播子相消，我们剩下
\begin{equation}
  \langle\mathbf{p}_1\mathbf{p}_2|S|\mathbf{k}_1\mathbf{k}_2\rangle
  = \text{所有切掉外腿的连通图之和},
  \label{eq:7.45}
\end{equation}
这正是主公式 \eqref{eq:4.103}。这就完成了 $S$ 矩阵元费曼规则的推导。

\section{光学定理}

$S$ 矩阵满足若干由幺正性推出的一般恒等式。其中最重要的是光学定理，它把前向散射振幅的虚部与总截面联系起来。在本节中，我们将证明这个定理，并展示它如何出现在费曼图展开中。

\subsection{幺正性与光学定理}

$S$ 矩阵是幺正的：$S^{\dagger}S = 1$。写出 $S = 1+iT$，这意味着
\begin{equation}
  -i(T - T^{\dagger}) = T^{\dagger}T .
\end{equation}
在初态 $|i\rangle$ 和末态 $|f\rangle$ 之间取矩阵元，并在右边插入一组完备态，我们得到
\begin{equation}
  -i\bigl( \langle f|T|i\rangle - \langle f|T^{\dagger}|i\rangle \bigr)
  = \sum_n \langle f|T^{\dagger}|n\rangle\langle n|T|i\rangle .
\end{equation}
对于前向散射振幅，$|f\rangle = |i\rangle$，这给出
\begin{equation}
  2\,\mathrm{Im}\,\mathcal{M}(i\to i)
  = \sum_n \int d\Pi_n\ |\mathcal{M}(i\to n)|^2 .
  \label{eq:7.49}
\end{equation}
利用 $|\mathcal{M}|^2$ 与截面之间的关系，我们得到光学定理：
\begin{equation}
  \mathrm{Im}\,\mathcal{M}(i\to i) = 2E_{\mathrm{cm}}\,p_{\mathrm{cm}}\,
    \sigma_{\mathrm{total}}(i\to\text{任意末态}),
  \label{eq:7.50}
\end{equation}
其中 $E_{\mathrm{cm}}$ 是总质心能量，$p_{\mathrm{cm}}$ 是质心系中任一粒子的动量。这个方程把前向散射振幅与所有末态产生的总截面联系起来。由于前向散射振幅的虚部给出波束穿过靶时前向波的衰减，这个量自然应当正比于散射概率。恒等式 \eqref{eq:7.50} 给出了精确的联系。

\subsection{费曼图的光学定理}

现在让我们研究 $S$ 矩阵元虚部的这个恒等式如何出现在费曼图展开中。很容易检验（例如在 QED 中），每个对 $S$ 矩阵元 $\mathcal{M}$ 有贡献的图都是纯实的，除非某些分母为零，从而使处理极点的 $i\epsilon$ 规则变得相关。因此，只有当图中的虚粒子在壳时，费曼图才对 $\mathcal{M}$ 的虚部有贡献。我们现在将展示如何分离并计算这个虚部。

为了我们当前的目的，让 $\mathcal{M}$ 由微扰论的费曼规则定义。这使我们能够把 $\mathcal{M}(s)$ 看成复变量 $s = E^2_{\mathrm{cm}}$ 的解析函数，尽管 $S$ 矩阵元只在外粒子具有实动量时才被定义。

我们首先证明，$\mathcal{M}(s)$ 虚部的出现总需要一个支割线奇点。设 $s_0$ 是最轻多粒子态产生的阈值能量。对于低于 $s_0$ 的实 $s$，中间态不能到壳，所以 $\mathcal{M}(s)$ 是实的。因此，对于实 $s < s_0$，我们有恒等式
\begin{equation}
  \mathcal{M}(s) = \bigl[ \mathcal{M}(s) \bigr]^{*} .
  \label{eq:7.51}
\end{equation}
这个方程的每一边都是 $s$ 的解析函数，所以它可以解析延拓到整个复 $s$ 平面。特别是，在 $s > s_0$ 的实轴附近，式~\eqref{eq:7.51} 蕴含
\begin{equation}
  \mathrm{Re}\,\mathcal{M}(s+i\epsilon) = \mathrm{Re}\,\mathcal{M}(s-i\epsilon);
  \qquad
  \mathrm{Im}\,\mathcal{M}(s+i\epsilon) = -\mathrm{Im}\,\mathcal{M}(s-i\epsilon).
\end{equation}
在实轴上有一条支割线，从阈值能量 $s_0$ 开始；割线两边的间断为
\begin{equation}
  \mathrm{Disc}\,\mathcal{M}(s) = 2i\,\mathrm{Im}\,\mathcal{M}(s+i\epsilon) .
\end{equation}
通常计算一个图的间断比直接计算虚部更容易。费曼传播子中的 $i\epsilon$ 规则表明，物理散射振幅应在割线上方、即 $s+i\epsilon$ 处求值。

我们在第7.1节已经看到，电子自能图具有从物理电子-光子阈值开始的支割线。现在让我们研究更一般的单圈图，并证明它们的间断恰好给出 \eqref{eq:7.49} 所要求的虚部。这个结果向多圈图的推广已由 Cutkosky 证明，%
\footnote{R.~E.~Cutkosky, \emph{J.\ Math.\ Phys.} \textbf{1}, 429 (1960).}他在证明过程中表明，费曼图在其支割线处的间断总是由一组简单的切割规则给出。%
\footnote{这些规则只对物理区域中的奇点是简单的。在物理区域之外，三点及更高点振幅的奇点会变得相当复杂。这个主题的综述见 R.~J.~Eden, P.~V.~Landshoff, D.~I.~Olive 和 J.~C.~Polkinghorne, \emph{The Analytic S-Matrix} (Cambridge University Press, 1966).}

我们首先在 $\phi^4$ 理论中检验 \eqref{eq:7.49}。由于 \eqref{eq:7.49} 的右边从 $\lambda^2$ 阶开始，我们预期 $\mathrm{Im}\,\mathcal{M}$ 的首次贡献也应来自高阶图。于是，考虑 $\lambda^2$ 阶图
\begin{equation}
  \delta\mathcal{M}(s) = \frac{(-i\lambda)^2}{2}\ \int
    \frac{d^4q}{(2\pi)^4}\ \frac{i}{(k/2-q)^2-m^2+i\epsilon}\
    \frac{i}{(k/2+q)^2-m^2+i\epsilon} ,
  \label{eq:7.52}
\end{equation}
其中圈在 $s$ 道中（$k = k_1+k_2$ 是总动量）。（很容易检验相应的 $t$ 道和 $u$ 道图在 $s$ 高于阈值时没有支割线奇点。）当用第6.3节的方法计算这个积分时，Wick 转动产生一个额外的 $i$ 因子，因此低于阈值时 $\delta\mathcal{M}$ 是纯实的。

我们想要验证积分 \eqref{eq:7.52} 在物理区域 $k^0 > 2m$ 内实轴上有一条间断。最容易的方法是先对 $k^0 < 2m$ 计算积分，然后通过解析延拓增大 $k^0$，从而识别这条间断。让我们在质心系中工作，其中 $k = (k^0, \mathbf{0})$。于是 \eqref{eq:7.52} 的被积函数在积分变量 $q^0$ 中有四个极点，位置在
\begin{equation}
  q^0 = \frac{1}{2}k^0 \pm (E_{\mathbf{q}} - i\epsilon),
  \qquad
  q^0 = -\frac{1}{2}k^0 \pm (E_{\mathbf{q}} - i\epsilon) .
\end{equation}
这些极点中有两个位于实 $q^0$ 轴上方，两个位于下方。我们将把积分围道向下闭合，并拾取下半平面中极点的留数。其中，只有 $q^0 = -\tfrac{1}{2}k^0 + E_{\mathbf{q}}$ 处的极点将对间断有贡献。注意，拾取这个极点的留数等价于在 $dq^0$ 积分下把
\begin{equation}
  \frac{1}{(k/2+q)^2-m^2+i\epsilon}
    \to -2\pi i\,\delta\bigl( (k/2+q)^2 - m^2 \bigr)
  \label{eq:7.53}
\end{equation}

这个极点的贡献给出积分
\begin{equation}
  \begin{aligned}
    i\delta\mathcal{M} &= -2\pi i\ \frac{\lambda^2}{2}\ \int
      \frac{d^3q}{(2\pi)^4}\ \frac{1}{2E_{\mathbf{q}}}\
      \frac{i}{(k^0-E_{\mathbf{q}})^2}\\
    &= -\frac{\lambda^2}{2}\ \frac{4\pi}{(2\pi)^4}\ \int_0^{\infty}
      dE_{\mathbf{q}}\ E_{\mathbf{q}}|\mathbf{q}|\ \frac{1}{2E_{\mathbf{q}}\,k^0\,
        (k^0-2E_{\mathbf{q}})} .
  \end{aligned}
  \label{eq:7.54}
\end{equation}
第二行中的被积函数在 $E_{\mathbf{q}} = k^0/2$ 处有一个极点。当 $k^0 < 2m$ 时，这个极点不在积分围道上，所以 $\delta\mathcal{M}$ 明显是实的。然而当 $k^0 > 2m$ 时，极点位于积分围道正上方或正下方，取决于 $k^0$ 被赋予小的正或负虚部。因此积分在 $k^2 + i\epsilon$ 与 $k^2 - i\epsilon$ 之间获得一条间断。为了计算这条间断，应用
\begin{equation}
  \frac{1}{k^0 - 2E_{\mathbf{q}} \pm i\epsilon}
  = \mathrm{P}\frac{1}{k^0 - 2E_{\mathbf{q}}}
    \mp i\pi\,\delta(k^0 - 2E_{\mathbf{q}})
\end{equation}
（其中 P 表示主值）。间断由把极点替换为 delta 函数给出。这反过来等价于把原来的传播子替换为 delta 函数：
\begin{equation}
  \frac{1}{(k/2-q)^2-m^2+i\epsilon}
    \to -2\pi i\,\delta\bigl( (k/2-q)^2 - m^2 \bigr) .
  \label{eq:7.55}
\end{equation}
现在让我们回顾一下我们所证明的步骤。回到原来的积分 \eqref{eq:7.52}，把两个传播子上的动量重新标为 $p_1$、$p_2$，并代换
\begin{equation}
  \frac{1}{p_i^2 - m^2 + i\epsilon}
    \to -2\pi i\,\delta(p_i^2 - m^2) .
  \label{eq:7.56}
\end{equation}
我们已经证明，积分的间断是通过把两个传播子各自替换为 delta 函数来计算的。$\mathcal{M}$ 的间断只来自 $d^4p_i$ 积分中两个 delta 函数同时满足的区域。通过对 delta 函数积分，我们使动量 $p_i$ 到壳，并把积分 $d^4p_i$ 转换为对相对论相空间的积分。表达式 \eqref{eq:7.52} 中剩下的是因子 $\lambda^2$，即领头阶散射振幅的平方，以及对称因子 $(1/2)$，它可以重新解释为适当的相空间归一化。这正是双体中间态的光学定理 \eqref{eq:7.49}。这就建立了单圈图的光学定理。

\section{沃德-高桥恒等式}

在本节中，我们将给出沃德恒等式 \eqref{eq:5.79} 的形式证明，实际上是一个更一般的关系——沃德-高桥恒等式的证明。考虑一个 QED 振幅 $\mathcal{M}$，它具有动量 $k$ 和极化 $\epsilon(k)$ 的外光子。写出 $\mathcal{M} = \epsilon^{*\mu}(k)\mathcal{M}_\mu(k)$，我们想证明
\begin{equation}
  k^\mu\mathcal{M}_\mu(k) = 0 .
  \label{eq:7.64}
\end{equation}
我们的证明逐图进行。首先考虑一条有 $n$ 个顶点贯穿而过的电子线，要给它附加一个额外外光子。电子传播子的动量是 $p, p_1 = p+q_1, p_2 = p_1+q_2$，依此类推直到 $p_f = p_{n-1}+q_n$。如果有 $n$ 个顶点，我们可以在 $n+1$ 个不同位置插入我们的光子。假设我们把它插在第 $i$ 个顶点之后。

现在让我们考察这些图的值，把极化矢量 $\epsilon^\mu(k)$ 替换为 $k^\mu$。$k^\mu$ 与新顶点的乘积可以方便地写成：
\begin{equation}
  -ie\,k^\mu\gamma_\mu = -ie\bigl[ (\slashed{p}_i+\slashed{k}-m)
    - (\slashed{p}_i-m) \bigr] .
  \label{eq:7.65}
\end{equation}
乘以相邻的电子传播子，我们得到
\begin{equation}
  \frac{i}{\slashed{p}_i+\slashed{k}-m}\ (-ie)\ \frac{i}{\slashed{p}_i-m}
  = ie\Bigl( \frac{1}{\slashed{p}_i-m}
    - \frac{1}{\slashed{p}_i+\slashed{k}-m} \Bigr) .
\end{equation}
因此，光子插在位置 $i$ 的图具有这样的结构：这个表达式的第一项与光子插在位置 $i-1$ 的图的第二项相消。任何其他一对具有相邻插入的图之间也会发生类似的相消。当我们对沿线的所有可能插入位置求和时，除端点处的未配对项外，一切都相消。来自最后一个顶点之后（最左边）插入的未配对项恰好是 $e$ 乘以原图的值；另一个未配对项来自第一个顶点之前的插入，除了一个负号以及处处把 $p$ 替换为 $p+k$ 之外完全相同。图解地，我们的结果是
\begin{equation}
  \sum_{\text{插入位置}} \text{（插入光子的图）}
  = e\Bigl[ \text{（动量为 }p\text{ 的原图）}
    - \text{（动量为 }p+k\text{ 的原图）} \Bigr].
  \label{eq:7.66}
\end{equation}
在 \eqref{eq:7.66} 左边的每个图中，进入电子线的动量是 $p$，离开的动量是 $q$。根据 LSZ 公式，我们可以通过取极点乘积
\begin{equation}
  \frac{1}{\slashed{p}-m}\ \frac{1}{\slashed{q}-m} .
\end{equation}
的系数，从每个图提取对 $S$ 矩阵元的贡献。\eqref{eq:7.66} 右边的项各自包含其中一个极点，但都不包含两个极点。因此 \eqref{eq:7.66} 的右边对 $S$ 矩阵没有贡献。%
\footnote{这一步论证只有在我们把微扰级数整理成传播子包含 $m$ 而非 $m_0$ 时才直截了当。我们将在第10章这样做。}

为了完成沃德-高桥恒等式的证明，我们必须考虑光子附加到内部电子圈的情形。在插入光子之前，一个典型圈有 $n$ 个顶点和 $n$ 个动量分别为 $p_1, p_1+q_2 = p_2$，依此类推直到 $p_n$ 的传播子。现在假设我们在顶点 $i$ 和 $i+1$ 之间插入光子 $\gamma(k)$。我们现在有一个额外的动量 $k$ 从新顶点绕圈流动。为了计算对圈的所有这种插入之和，对每个图应用恒等式 \eqref{eq:7.65}。一个图的第一项将被光子插在下一个顶点处的图的项之一相消。来自其他相邻插入对的项之间也会发生类似的相消。当我们对全部 $n$ 个插入位置求和时，剩下
\begin{equation}
  \text{（光子插在顶点 }n\text{ 和顶点 }1\text{ 之间的图）}
  + \text{（光子插在顶点 }1\text{ 和顶点 }2\text{ 之间的图）},
  \label{eq:7.67}
\end{equation}
它们只相差圈动量移动 $k$。在第二项中把积分变量从 $p_1$ 移到 $p_1+k$，我们看到剩下的两项相消。因此，光子沿闭合圈插入的图之和为零。

我们现在准备把证明的各部分组装起来。假设振幅 $\mathcal{M}$ 有 $2n$ 条外电子线，$n$ 条入射和 $n$ 条出射。把入射动量标为 $p_i$，出射动量标为 $q_i$：
\begin{equation}
  \mathcal{M}_\mu(k;\ p_1\cdots p_n;\ q_1\cdots q_n) .
\end{equation}
（振幅还可以涉及任意数目的额外外光子。）振幅 $\mathcal{M}_0$ 缺少光子 $\gamma(k)$，但在其他方面完全相同。为了由 $\mathcal{M}_0$ 构成 $k^\mu\mathcal{M}_\mu$，我们必须对所有对 $\mathcal{M}_0$ 有贡献的图求和，并且对每个图，对光子可能插入的所有位置求和。对任何图中内圈上的插入位置求和给出零。对任何图中贯穿而过的线上的插入位置求和给出 \eqref{eq:7.66} 形式的贡献。对任何特定图的所有插入位置求和，我们得到
\begin{equation}
  \text{（一个外动量移动了的图之差）} .
\end{equation}
对所有这样的图求和，我们最终得到
\begin{equation}
  \begin{aligned}
    k_\mu\mathcal{M}^{\mu}(k;\ p_1\cdots p_n;\ q_1\cdots q_n)
    &= e\ \sum_{i} \Bigl[ \mathcal{M}_0(p_1\cdots p_n;\ q_1\cdots
      (q_i-k)\cdots)\\
    &\quad - \mathcal{M}_0((p_i+k)\cdots p_n;\ q_1\cdots q_n) \Bigr].
  \end{aligned}
  \label{eq:7.68}
\end{equation}
这就是 QED 中关联函数的沃德-高桥恒等式。我们在 \eqref{eq:7.66} 下面看到右边对 $S$ 矩阵没有贡献；因此在 $\mathcal{M}$ 是 $S$ 矩阵元的特殊情形下，式~\eqref{eq:7.68} 简化为沃德恒等式 \eqref{eq:7.64}。

在进一步讨论这个恒等式之前，我们应该提及上述证明中的一个潜在缺陷。为了在式 \eqref{eq:7.67} 中找到所需的相消，我们必须把积分变量平移一个常数。然而，如果积分发散，这种平移是不允许的。类似地，在导致式~\eqref{eq:7.66} 的表达式中可能存在发散的圈动量积分。这里在证明中没有显式的平移，但在实践中，人们在计算积分时通常确实会作平移。无论哪种情况，紫外发散都可能使沃德-高桥恒等式失效。我们将在下一节看到这个问题的例子，以及它的一般解决方案。

\section{电荷重整化}

本章的最后一个主题是真空极化图的计算，即 \eqref{eq:6.1} 中的第四个图。这个图也被称为对光子传播子的修正。设 $-i\Pi^{\mu\nu}(q)$ 表示对光子传播子的所有 1PI 插入之和：
\begin{equation}
  -i\Pi^{\mu\nu}(q) = \text{（所有 1PI 光子自能图之和）} .
  \label{eq:7.71}
\end{equation}
洛伦兹不变性和流守恒约束 $\Pi^{\mu\nu}(q)$ 的形式：
\begin{equation}
  \Pi^{\mu\nu}(q) = (q^2g^{\mu\nu} - q^\mu q^\nu)\,\Pi(q^2) .
  \label{eq:7.72}
\end{equation}
（任何不含补偿项 $q^\mu q^\nu$ 的正比于 $g^{\mu\nu}$ 的项都会违反沃德恒等式并给光子一个质量。）精确光子传播子通过求和 1PI 插入的几何级数得到：
\begin{equation}
  \frac{-ig_{\mu\nu}}{q^2}\ \Bigl[ 1 + \Pi(q^2) + \Pi(q^2)^2 + \cdots \Bigr]
  = \frac{-ig_{\mu\nu}}{q^2\,[1 - \Pi(q^2)]} .
  \label{eq:7.73}
\end{equation}
在任何 $S$ 矩阵元计算中，这个精确传播子至少有一个端点会连接到费米子线。当我们对线上它可能连接的所有位置求和时，根据沃德恒等式，我们必定发现正比于 $q_\mu$ 或 $q_\nu$ 的项消失。因此，为了计算 $S$ 矩阵元的目的，我们可以缩写
\begin{equation}
  \frac{-ig_{\mu\nu}}{q^2\,[1-\Pi(q^2)]}
  \xrightarrow{\text{沃德恒等式}}
  \frac{-ig_{\mu\nu}}{q^2}\ \frac{1}{1-\Pi(q^2)} .
  \label{eq:7.75}
\end{equation}
注意，只要 $\Pi(q^2)$ 在 $q^2 = 0$ 处正则，精确传播子就总有一个 $q^2 = 0$ 的极点。换句话说，光子在所有微扰论阶都保持严格无质量。这个论断强烈依赖于我们在 \eqref{eq:7.73} 中对沃德恒等式的使用。例如，如果 $\Pi^{\mu\nu}(q)$ 包含一项 $M^2g^{\mu\nu}$（没有补偿的 $q^\mu q^\nu$ 项），光子质量将被移动到 $M$。$q^2 = 0$ 极点的留数是
\begin{equation}
  \frac{1}{1 - \Pi(0)} = Z_3 .
  \label{eq:7.75b}
\end{equation}
任何低 $q^2$ 散射过程的振幅相对于树图近似都将被这个因子移动：
\begin{equation}
  e_0^2 \to e_0^2\,Z_3 .
\end{equation}
由于在光子传播子的每个端点都有一个 $e$ 因子，我们可以方便地通过作替换 $e_0 \to \sqrt{Z_3}\,e_0$ 来计入这个移动。这个替换被称为\emph{电荷重整化}；它在许多方面类似于第7.1节引入的质量重整化。特别要注意，实验中测得的“物理”电子电荷是 $\sqrt{Z_3}\,e_0$。因此我们将移动记号，把这一量直接称为 $e$。从现在起，我们将把“裸”电荷（拉格朗日量中乘以 $\gamma^\mu$ 的量）称为 $e_0$。于是我们有
\begin{equation}
  (\text{物理电荷}) = e = \sqrt{Z_3}\,e_0 = \sqrt{Z_3}\,
    (\text{裸电荷}) .
  \label{eq:7.76}
\end{equation}
在最低阶，$Z_3 = 1$ 且 $e = e_0$。

除了这个电子电荷强度的常数移动之外，$\Pi(q^2)$ 还有另一个效应。考虑一个具有非零 $q^2$ 的散射过程，并假设我们已经把 $\Pi(q^2)$ 计算到 $\alpha$ 的领头阶：$\Pi(q^2) \approx \Pi_2(q^2)$。该过程的振幅将涉及量
\begin{equation}
  \frac{e_0^2}{q^2}\ \frac{1}{1-\Pi(q^2)} = \frac{e^2}{q^2}\
    \frac{1}{1 - [\Pi_2(q^2) - \Pi_2(0)]} .
\end{equation}
（在最低阶，把 $e_0^2$ 换成 $e^2$ 无关紧要。）括号中的量可以解释为一个依赖 $q^2$ 的电荷。因此，用精确光子传播子替换树图光子传播子的全部效应就是替换
\begin{equation}
  \alpha = \frac{e_0^2}{4\pi}
    \to \alpha_{\mathrm{eff}}(q^2) = \frac{e^2}{4\pi}\
      \frac{1}{1 - [\Pi_2(q^2)-\Pi_2(0)]}
    = \frac{\alpha}{1 - [\Pi_2(q^2)-\Pi_2(0)]} .
  \label{eq:7.77}
\end{equation}
（在领头阶，我们同样可以把 $\Pi$ 项移到分子；但我们将在第12章看到，在这种形式下，当把 $\Pi_2$ 替换为 $\Pi$ 时，该表达式对所有阶都成立。）

\subsection{$\Pi_2$ 的计算}

既然我们如此费力地解释了 $\Pi_2(q^2)$，最好还是把它计算出来。回到 \eqref{eq:7.71}，我们有
\begin{equation}
  \begin{aligned}
    -i\Pi^{\mu\nu}_2(q) &= -(-ie)^2 \int \frac{d^4k}{(2\pi)^4}\
      \mathrm{tr}\Bigl[ \gamma^\mu\ \frac{i(\slashed{k}+m)}{k^2-m^2}\
        \gamma^\nu\ \frac{i(\slashed{k}+\slashed{q}+m)}
          {(k+q)^2-m^2} \Bigr]\\
    &= -4e^2 \int \frac{d^4k}{(2\pi)^4}\
      \frac{k^\mu(k+q)^\nu + k^\nu(k+q)^\mu - g^{\mu\nu}
        (k\cdot(k+q) - m^2)}
        {(k^2-m^2)((k+q)^2-m^2)} .
  \end{aligned}
  \label{eq:7.78}
\end{equation}
为了方便，我们写 $e$ 和 $m$ 而不是 $e_0$ 和 $m_0$，因为差别只给 $\Pi_2$ 带来 $\alpha^2$ 阶的贡献。

现在引入一个费曼参数来合并分母因子：
\begin{equation}
  \frac{1}{(k^2-m^2)((k+q)^2-m^2)}
  = \int_0^1 dx\ \frac{1}{(k^2 + 2xk\cdot q + xq^2 - m^2)^2}
  = \int_0^1 dx\ \frac{1}{(\ell^2 + x(1-x)q^2 - m^2)^2},
\end{equation}
其中 $\ell = k + xq$。用 $\ell$ 表示，\eqref{eq:7.78} 的分子为
\begin{equation}
  \text{分子} = 2\ell^\mu\ell^\nu - 2x(1-x)q^\mu q^\nu
    + g^{\mu\nu}\bigl( m^2 + x(1-x)q^2 \bigr)
    + (\text{线性于 } \ell \text{ 的项}) .
\end{equation}
作 Wick 转动并代换 $\ell_E^0 = -i\ell^0$，我们得到
\begin{equation}
  -i\Pi^{\mu\nu}_2(q) = 4e^2 \int_0^1 dx\ \int
    \frac{d^4\ell_E}{(2\pi)^4}\ \frac{1}{(\ell_E^2+\Delta)^2}\
    \Bigl[ 2\ell_E^\mu\ell_E^\nu - g^{\mu\nu}\ell_E^2 + \cdots \Bigr],
  \label{eq:7.79}
\end{equation}
其中 $\Delta = m^2 - x(1-x)q^2$。这个积分的紫外发散非常严重。如果我们把它截断在 $\ell_E = \Lambda$，我们会发现领头项为
\begin{equation}
  i\Pi^{\mu\nu}_2(q) \propto e^2\Lambda^2 g^{\mu\nu},
\end{equation}
没有补偿的 $q^\mu q^\nu$ 项。这个结果违反沃德恒等式；它将给光子一个无穷大质量 $M \propto e\Lambda$。

我们对沃德恒等式的证明失败了，其方式正如同上一节末尾所预料的那样：当积分发散时，\eqref{eq:7.67} 中积分变量的平移是不允许的。在我们现在的计算中，沃德恒等式的失效是灾难性的：它导致无穷大光子质量，%
\footnote{我们仍然可以通过在拉格朗日量中加入一个补偿的无穷大光子质量项来使观测到的光子质量为零。更一般地，我们可以向拉格朗日量加入一些项，使沃德恒等式对任何 $n$ 点关联函数都成立。这一程序会给出与我们即将采用的方法相同的结果，但要复杂得多。}这与实验相冲突。幸运的是，有一种方法可以挽救沃德恒等式。

在上面的分析中，我们以最直接也最天真的方式对发散积分进行了正则化：把它截断在很大的动量 $\Lambda$ 处。但其他正则子也是可能的，而且有些实际上会保持沃德恒等式。在我们对顶点图和电子自能图的计算中，我们使用了 Pauli--Villars 正则子。这个正则子保持了 $Z_1 = Z_2$ 这一关系，它是沃德恒等式的一个推论；而天真的截断则不保持（见习题7.2）。我们可以通过引入 Pauli--Villars 费米子来修正我们现在的计算。不幸的是，这需要若干组费米子，使方法相当复杂。%
\footnote{这一方法见 Bjorken and Drell (1964), p.~154.}我们将使用一种更简单的方法，即归功于 't Hooft 和 Veltman 的\emph{维数正则化}。%
\footnote{G.~'t~Hooft and M.~J.~G.~Veltman, \emph{Nucl.\ Phys.} \textbf{B44}, 189 (1972).}维数正则化保持 QED 以及一大类更一般理论的对称性。

在量子场论中，使用哪个正则子的问题没有先验答案。通常这个选择对理论的预言没有影响。当两个正则子对可观测量给出不同答案时，通常是因为某个对称性（如沃德恒等式）被其中一个（或两个）破坏了。在这些情形下，我们把这个对称性视为基本的，并要求正则子保持它。但这个选择的正确性无法证明；我们是在把这个对称性作为一条新公理来采用。

\subsection{维数正则化}

维数正则化的思想非常容易陈述：把费曼图作为时空维数 $d$ 的解析函数来计算。对于足够小的 $d$，任何圈动量积分都会收敛，因此沃德恒等式可以被证明。任何可观测量的最终表达式在 $d\to 4$ 时都应当具有良好定义的极限。

让我们做一个练习计算，看看这个技术如何运作。我们认为时空有一个时间维和 $(d-1)$ 个空间维。然后我们可以像之前一样对费曼积分作 Wick 转动，给出对 $d$ 维欧几里得空间的积分。一个典型的例子是
\begin{equation}
  \int \frac{d^d\ell_E}{(2\pi)^d}\ \frac{1}{(\ell_E^2 + \Delta)^2}
  = \int_0^{\infty} \frac{d\ell_E}{(2\pi)^d}\ \frac{S_d\,\ell_E^{d-1}}
    {(\ell_E^2 + \Delta)^2} ,
  \label{eq:7.80}
\end{equation}
其中 $S_d = 2\pi^{d/2}/\Gamma(d/2)$ 是 $d$ 维中单位球面的面积。为了计算它，使用下面的技巧：
\begin{equation}
  (\sqrt{\pi})^d = \Bigl( \int dx\ e^{-x^2} \Bigr)^d
    = \int d^d x\ e^{-|x|^2}
    = S_d \int_0^{\infty} dr\ r^{d-1} e^{-r^2}
    = S_d\cdot \tfrac{1}{2}\Gamma(d/2) .
\end{equation}
利用标准公式
\begin{equation}
  \int_0^{\infty} d\ell\ \frac{\ell^{d-1}}{(\ell^2+\Delta)^2}
    = \Delta^{d/2-2}\ \frac{\Gamma(d/2)\Gamma(2-d/2)}{2},
\end{equation}
计算对 $\ell_E$ 的积分，我们得到
\begin{equation}
  \int \frac{d^d\ell_E}{(2\pi)^d}\ \frac{1}{(\ell_E^2+\Delta)^2}
  = \frac{\Gamma(2-d/2)}{(4\pi)^{d/2}}\ \Delta^{d/2-2} .
  \label{eq:7.81}
\end{equation}
当 $d\to 4$ 时，伽马函数 $\Gamma(2-d/2)$ 有一个极点。令 $d = 4-2\epsilon$ 在 $d=4$ 附近展开，我们利用
\begin{equation}
  \Gamma(\epsilon) = \frac{1}{\epsilon} - \gamma_E + \mathcal{O}(\epsilon),
\end{equation}
其中 $\gamma_E = 0.5772\ldots$ 是 Euler--Mascheroni 常数。发散表现为 $\epsilon$ 中的极点，它将被重整化吸收。

把维数正则化应用于真空极化积分 \eqref{eq:7.79}，并收缩以 $(q^2g^{\mu\nu} - q^\mu q^\nu)$，我们得到
\begin{equation}
  \Pi_2(q^2) = -\frac{2\alpha}{\pi}\ \int_0^1 dx\
    x(1-x)\ \Bigl[ \frac{1}{\epsilon} - \gamma_E + \log\Bigl(
      \frac{4\pi\mu^2}{\Delta} \Bigr) + \cdots \Bigr],
  \label{eq:7.82}
\end{equation}
其中 $\mu$ 是由维数正则化引入的质量标度。物理的、重整化的结果通过减去 $q^2 = 0$ 处的值得到：
\begin{equation}
  \Pi_2(q^2) - \Pi_2(0) = -\frac{2\alpha}{\pi}\ \int_0^1 dx\
    x(1-x)\ \log\Bigl( \frac{m^2}{m^2 - x(1-x)q^2} \Bigr) .
  \label{eq:7.91}
\end{equation}
完成积分，
\begin{equation}
  \Pi_2(q^2) - \Pi_2(0) = -\frac{2\alpha}{\pi}\ \Bigl[
    \frac{5}{9} - \frac{4m^2}{3q^2} + \frac{(4m^2/q^2)(1+2m^2/q^2)}
      {\sqrt{1-4m^2/q^2}}\ \mathrm{arctanh}\Bigl(
        \frac{1}{\sqrt{1-4m^2/q^2}} \Bigr) \Bigr],
  \label{eq:7.91b}
\end{equation}
对 $q^2 < 0$ 成立，并解析延拓到物理区域。

让我们计算 $\Pi_2$ 在 $q^2 > 4m^2$ 时的虚部。对于任何固定的 $q^2$，有贡献的 $x$ 值位于 $x = \tfrac{1}{2}\pm\tfrac{1}{2}\beta$ 两点之间，其中
$\beta = \sqrt{1-4m^2/q^2}$。由于
$\mathrm{Im}\,[\log(-X\pm i\epsilon)] = \pm\pi$，我们有
\begin{equation}
  \begin{aligned}
    \mathrm{Im}\,[\Pi_2(q^2\pm i\epsilon)]
    &= \mp \frac{2\alpha}{\pi}\cdot\pi\ \int_{\frac{1}{2}-\frac{\beta}{2}}
      ^{\frac{1}{2}+\frac{\beta}{2}} dx\ x(1-x)\\
    &= \mp 2\alpha\ \int_{-\beta/2}^{\beta/2} dy\ \Bigl(
      \frac{1}{4} - y^2 \Bigr)
    = \mp \frac{2\alpha}{3}\ \beta\Bigl( 1 - \frac{\beta^2}{4} \Bigr)\ \cdots
    = \mp \frac{\alpha}{3}\ \beta(1 - \tfrac{1}{3}\beta^2)\ \cdots .
  \end{aligned}
  \label{eq:7.92}
\end{equation}
这个对 $q^2$ 的依赖与式~\eqref{eq:5.13}——费米子-反费米子对产生的截面——中的依赖完全相同。这正是我们从图~\ref{fig:7.6}(b) 所示的幺正关系中所预期的；穿过前向 Bhabha 散射图的割线给出 $e^{+}e^{-}\to f\bar{f}$ 的总截面。参数 $\beta$ 恰好是质心系中费米子的速度。

\begin{figure}[t]
  \centering
  \includegraphics[width=0.6\textwidth]{fig7.6.png}
  \caption{Bhabha 散射光学定理的两个贡献。}
  \label{fig:7.6}
\end{figure}

接下来让我们考察 $\Pi_2(q^2)$ 如何修改电磁相互作用，如式~\eqref{eq:7.77} 所确定的。在非相对论极限下，计算势 $V(\mathbf{r})$ 是有意义的。对于异号电荷之间的相互作用，与式~\eqref{eq:4.135} 类比，我们有
\begin{equation}
  V(\mathbf{x}) = \int \frac{d^3q}{(2\pi)^3}\ e^{i\mathbf{q}\cdot\mathbf{x}}\
    \frac{-e^2}{|\mathbf{q}|^2\,[1 - \Pi_2(-|\mathbf{q}|^2)]} .
  \label{eq:7.93}
\end{equation}
对 $|q^2| \ll m^2$ 展开 $\Pi_2$，我们得到
\begin{equation}
  V(\mathbf{r}) = -\frac{\alpha}{r}\ \Bigl[ 1 + \frac{4\alpha}{5\pi}\,
    \frac{e^{-2mr}}{(mr)^{3/2}}\ \cdots \Bigr]
  \approx -\frac{\alpha}{r}\ -\ \frac{4\alpha^2}{5\sqrt{\pi}m^{3/2}}\,
    \delta^{(3)}(\mathbf{x}) .
  \label{eq:7.94}
\end{equation}
修正项表明电磁力在小距离处变得强得多。这个效应可以在氢原子中测量，其中能级被移动
\begin{equation}
  \Delta E = \int d^3x\ |\psi(\mathbf{x})|^2\ \Bigl(
    -\frac{4\alpha^2}{5\sqrt{\pi}m^{3/2}}\delta^{(3)}(\mathbf{x}) \Bigr)
  = -\frac{4\alpha^2}{5\sqrt{\pi}m^{3/2}}\ |\psi(0)|^2 .
\end{equation}
波函数 $\psi(\mathbf{x})$ 只有在 $s$ 波态才在原点处非零。对于 $2S$ 态，移动为
\begin{equation}
  \Delta E = -\frac{4\alpha^2}{5\sqrt{\pi}m^{3/2}}\ \frac{1}{\pi a_0^3}\ \cdots
  = -\frac{\alpha^5 m}{15\sqrt{2}\pi}\ \cdots .
\end{equation}
这是表~\ref{tab:6.1} 所列 Lamb 位移劈裂的一（小）部分。

\begin{figure}[t]
  \centering
  \includegraphics[width=0.5\textwidth]{fig7.7.png}
  \caption{在非相对论极限下计算电磁相互作用有效强度的围道。$Q = i\mu$ 处的极点给出库仑势。支割线给出由真空极化引起的 $\alpha$ 阶修正。}
  \label{fig:7.7}
\end{figure}

式~\eqref{eq:7.94} 中的 delta 函数只是一个近似；为了找到修正项的真实范围，我们把式~\eqref{eq:7.93} 写成形式
\begin{equation}
  V(\mathbf{r}) = \int \frac{d^3q}{(2\pi)^3}\ e^{i\mathbf{q}\cdot\mathbf{r}}\
    \frac{-\alpha}{|\mathbf{q}|^2 + \mu^2}\ \Bigl[ 1 +
      \Pi_2(-|\mathbf{q}|^2) + \cdots \Bigr],
\end{equation}
其中我们插入了光子质量 $\mu$ 来对库仑势正则化。为了完成这个积分，我们把围道向上推（见图~\ref{fig:7.7}）。领头贡献来自 $Q = i\mu$ 处的极点，给出库仑势 $-\alpha/r$。但支割线还有额外贡献，它起始于 $Q = 2mi$。被积函数的实部在割线两侧相同，因此对积分的唯一贡献来自 $\Pi_2$ 的虚部。定义 $q = -iQ$，我们发现来自割线的贡献为
\begin{equation}
  \begin{aligned}
    \delta V(\mathbf{r}) &= \frac{-\alpha}{\pi^2 r}\ \int_{2m}^{\infty} dq\
      \frac{\mathrm{Im}\,\Pi_2(-q^2)}{q^2}\ \sin(qr)\ \cdots
    = \frac{-\alpha}{\pi r}\ \int_{2m}^{\infty} dq\
      \frac{2\alpha}{3}\ \frac{\beta(1-\tfrac{1}{3}\beta^2)}{q^2}\ \sin(qr)\ \cdots .
  \end{aligned}
\end{equation}
当 $r \ll 1/m$ 时，这个积分由 $q\sim 2m$ 附近的区域主导。在这个区域近似被积函数并代换 $t = q - 2m$，我们得到
\begin{equation}
  \delta V(\mathbf{r}) = \frac{-\alpha}{r}\ \frac{4\alpha}{\sqrt{\pi}}\
    \int_0^{\infty} dt\ \frac{e^{-(t+2m)r}}{(mr)^{3/2}}\ \cdots
  = -\frac{\alpha}{r}\ \frac{4\alpha}{5\sqrt{\pi}}\,
    \frac{e^{-2mr}}{(mr)^{3/2}} ,
  \label{eq:7.95}
\end{equation}
因此
\begin{equation}
  V(\mathbf{r}) = -\frac{\alpha}{r}\ \Bigl[ 1 + \frac{4\alpha}{5\sqrt{\pi}}\,
    \frac{e^{-2mr}}{(mr)^{3/2}} \Bigr] .
  \label{eq:7.95b}
\end{equation}
因此修正项的范围大约是电子康普顿波长 $1/m$。由于氢原子波函数在这个尺度上几乎不变，式~\eqref{eq:7.94} 中的 delta 函数是一个很好的近似。对 $V(\mathbf{r})$ 的辐射修正被称为\emph{乌林势}（Uehling potential）。

\begin{figure}[t]
  \centering
  \includegraphics[width=0.5\textwidth]{fig7.8.png}
  \caption{虚 $e^{+}e^{-}$ 对实际上是长度 $\sim 1/m$ 的偶极子，它们屏蔽电子的裸电荷。}
  \label{fig:7.8}
\end{figure}

我们可以把这个修正解释为屏蔽效应。在 $r\gg 1/m$ 处，虚 $e^{+}e^{-}$ 对使真空成为介电介质，其中表观电荷小于真实电荷（见图~\ref{fig:7.8}）。在更小的距离处，我们开始穿透极化云并看到裸电荷。这个现象被称为\emph{真空极化}。

现在考虑相反的极限：小距离或 $-q^2 \gg m^2$。式 \eqref{eq:7.91} 于是变成
\begin{equation}
  \Pi_2(q^2) \approx \frac{2\alpha}{\pi}\ \int_0^1 dx\
    x(1-x)\ \Bigl[ \log\Bigl( \frac{-q^2}{m^2} \Bigr)
      + \log(x(1-x)) + \mathcal{O}\Bigl( \frac{m^2}{q^2} \Bigr) \Bigr] .
\end{equation}
因此这个极限下的有效耦合常数为
\begin{equation}
  \alpha_{\mathrm{eff}}(q^2) = \frac{\alpha}{1 +
    \frac{2\alpha}{3\pi}\log\Bigl( \frac{-q^2}{m^2} \Bigr)
      + \cdots }
  \xrightarrow{-q^2 \gg m^2}
    \frac{\alpha}{1 - \frac{\alpha}{3\pi}\log\Bigl(
      \frac{-q^2}{\Lambda^2} \Bigr)} ,
  \label{eq:7.96}
\end{equation}
其中 $\Lambda = \exp(5/3)\,m$。有效电荷在小距离处变得大得多，因为我们穿透了虚电子-正电子对的屏蔽云。

\begin{figure}[t]
  \centering
  \includegraphics[width=0.7\textwidth]{fig7.9.png}
  \caption{HRS 合作组（M.~Derrick 等人，\emph{Phys.\ Rev.\ D} \textbf{34}，3286 (1986)）在 $E_{\mathrm{cm}} = 29$ GeV 处测量的 Bhabha 散射 $e^{+}e^{-}\to e^{+}e^{-}$ 的微分截面。上面的曲线是习题5.2中导出的 $\alpha^2$ 阶预言，加上弱相互作用的很小（$\sim 2\%$）修正。下面的曲线包括除真空极化贡献之外的所有 $\alpha^3$ 阶 QED 辐射修正；注意这些修正依赖于实验条件，如第6章所解释。中间的曲线还包括真空极化贡献，它在这个能量下把 $\alpha^2$ 的有效值增大约 $10\%$。}
  \label{fig:7.9}
\end{figure}

电子以及更重夸克和轻子的真空极化联合效应使 $\alpha_{\mathrm{eff}}(q^2)$ 的值从 $q = 0$ 到 $q = 30$ GeV 增大约 $5\%$，这个效应在高能实验中被观测到。图~\ref{fig:7.9} 显示了 $E_{\mathrm{cm}} = 29$ GeV 处 Bhabha 散射的截面，以及与包含和不包含真空极化图的 QED 的比较。我们可以通过令 $q = 1/r$ 把 $\alpha_{\mathrm{eff}}$ 写成 $r$ 的函数。$\alpha_{\mathrm{eff}}(r)$ 对所有 $r$ 的行为画在图~\ref{fig:7.10} 中。依赖距离（或“依赖标度”或“跑动”）的耦合常数的思想将是本书其余部分的一个主要主题。

\begin{figure}[t]
  \centering
  \includegraphics[width=0.6\textwidth]{fig7.10.png}
  \caption{由单圈真空极化图产生的有效电磁耦合常数随距离变化的定性草图。水平标度覆盖许多数量级。}
  \label{fig:7.10}
\end{figure}

\problems

\textbf{7.1} 在第7.3节中，我们用一种间接方法分析了 $\phi^4$ 理论中玻色子-玻色子散射的单圈 $s$ 道图。为了验证我们间接分析的正确性，用标准的费曼参数方法计算全部三个单圈图。检验光学定理的有效性。

\textbf{7.2} \emph{QED 中的替代正则子。} 在第7.5节中，我们看到沃德恒等式可能被不恰当选择的正则子破坏。让我们对若干正则子的选择，在 $\alpha$ 阶检验恒等式 $Z_1 = Z_2$ 的有效性。我们已经验证了该关系对 Pauli--Villars 正则化成立。
\begin{enumerate}[label=(\alph*)]
  \item 重新计算 $\delta Z_1$ 和 $\delta Z_2$，通过简单地对 $|\ell_E|$ 的积分设置上限 $\Lambda$ 来定义式 (6.49) 和式 (6.50) 的积分。证明在这个定义下 $\delta Z_1 \neq \delta Z_2$。
  \item 重新计算 $\delta Z_1$ 和 $\delta Z_2$，用维数正则化定义式 (6.49) 和式 (6.50) 的积分。你可以像往常一样取 Dirac 矩阵为 $4\times4$，但注意在 $d$ 维中，
    \begin{equation}
      g^\mu_{\ \mu}\gamma^\nu\gamma_\nu = d .
    \end{equation}
    证明在这个定义下 $\delta Z_1 = \delta Z_2$。
\end{enumerate}

\textbf{7.3} 在正文中我们讨论了涉及一个外光子的振幅的沃德恒等式。有一个重要的推广到涉及\emph{两个}外光子的振幅，称为光子-光子振幅的沃德-高桥恒等式。考虑两个动量分别为 $k_1$ 和 $k_2$ 的光子产生的振幅 $\mathcal{M}^{\mu\nu}(k_1, k_2)$。证明该振幅满足
\begin{equation}
  k_{1\mu}\mathcal{M}^{\mu\nu}(k_1, k_2) = 0 .
\end{equation}
（你可能想利用正文中的结果：沿贯穿而过的费米子线的光子插入之和是望远镜式的，而沿闭合圈的插入之和为零。）

\textbf{7.4} 验证正文中真空极化的结果，式~\eqref{eq:7.91}，在物理区域 $q^2 > 4m^2$ 内满足幺正关系
\begin{equation}
  \mathrm{Im}\,\Pi_2(q^2) = \frac{1}{2}\ \int d\Pi_2\
    \mathcal{M}(e^{+}e^{-}\to f\bar{f})\ ,
\end{equation}

\textbf{7.5} 在本习题中，我们将推导第7.1节所引用的场强重整化常数 $Z$ 与物理传播子极点之间的关系。
\begin{enumerate}[label=(\alph*)]
  \item 考虑精确的两点函数 $D(p^2)$，写出它的 K\"all\'en--Lehmann 谱表示。
  \item 通过在极点附近展开谱表示，证明 $p^2 = m^2$ 处极点的留数是 $iZ$。
  \item 证明重整化的传播子，定义为 $D_R(p^2) = Z^{-1}D(p^2)$，在极点处具有单位留数。
\end{enumerate}
